\documentclass[12pt, a4paper]{book}
\usepackage[english, german]{babel}
%\usepackage[latin1]{inputenc}
\usepackage[T1]{fontenc}
\usepackage{pdfpages}
\usepackage[bookmarksopenlevel=1]{hyperref}
\usepackage[margin=1in]{geometry}
\usepackage[numbered]{bookmark}
\usepackage{titlesec}
\usepackage{graphicx}

\newcommand{\thesubject}{Verantstaltung}

% LuaTex Specifics
\usepackage[math]{fontspec}
\usepackage{unicode-math}
\setmathfont{Asana-Math.otf}
\usepackage{fontenc}
\setmainfont{DINOT-Light.otf}[
	BoldFont = DINOT-Medium.otf,
	ItalicFont = DINOT-Regular.otf]
	
\titleformat{\chapter}[display]
  {\vskip-30pt\fontspec{DINOT-Regular.otf}\large}
  {\parbox[b]{0.5\linewidth}{\thesubject}\hfill \includegraphics[scale=1]{logos/Logo_HSRT_Schwarz}\vspace{40pt} \huge \chaptername~\thechapter}{2pt}{\fontspec{DINOT-Medium.otf}\Huge}
\titleformat{\section}
  {\fontspec{BlenderPro-Medium.otf}\Large}
  {\thesection}{1em}{}


\hypersetup{
  pdftitle    = {Skript zur Vorlesung},
  pdfsubject  = {Regelungstechnik und -systeme},
  pdfauthor   = {Prof. A. Notholt},
  pdfkeywords = {} ,
  pdfcreator  = {pdflatex},
  pdfproducer = {Reutlingen-University.de}
}


\begin{document}

	\renewcommand{\thesubject}{Regelungstechnik 1\\ Regelungstechnik 2 \\Regelungssysteme}
	
	%%%%% EINFÜHRUNG %%%%	
	
\chapter{Einführung in die Regelungstechnik}

	Dieses Kapitel legt die Grundsteine für einen gemeinsamen Sprache an. In diesem Kapitel werden die wichtigste Konzepte und Begriffe erklärt. Das Ziel ist, dass alle Studierende das gleiche darunter verstehen, und dies über die ganze Vorlesungsreihe.
	
\section{Was \glqq Regelungstechnik\grqq~ ist}

Die Regelungstechnik ist ein Fachbereich, der sich mit der Aufgabe befasst, \glqq einen zeitlich verändernden Prozess so zu beeinflussen, dass dieser Prozess in einer vorgegebenen Weise abläuft\grqq~\cite{LunzeRT2014}. Typische regelungstechnische Aufgaben findet man überall, wo ein gewünschtes Verhalten erreicht werden soll.

	
	%%% A VER %%%%
	
\end{document}
